\section{Formulaciones del modelo (F1--F4)}
\label{sec:formulaciones}

Presentamos cuatro formulaciones MILP equivalentes en óptimo entero pero con \emph{estructura}
muy distinta. La formulación final usada como base de la descomposición es \textbf{F3}; su
relajación, escrita de forma compacta sobre $\theta$, es \textbf{F4}.

\paragraph{Marco del curso: relajación LP, brecha de integridad y fortaleza.} Para un MILP de
minimización, reemplazar $y_j\in\{0,1\}$ por $y_j\in[0,1]$ da la \emph{relajación lineal}, cuya
solución $v(\text{LP})$ es una \textbf{cota inferior}: $v(\text{LP})\le v(\text{MILP})$. La
\textbf{brecha de integridad} $v(\text{MILP})-v(\text{LP})\ge0$ mide la fortaleza de la
formulación: brecha pequeña $\Rightarrow$ formulación \emph{fuerte} $\Rightarrow$ mejores cotas y
menos nodos en branch-and-bound. Un hecho central de esta sección es que \textbf{F1, F2, F3 y F4
tienen exactamente la misma relajación LP} (y por tanto la misma brecha de integridad); difieren
solo en cuántas variables y qué tan densa es su matriz, lo que el solver explota de forma muy
distinta en la práctica.

\subsection{F1 — formulación clásica (ReVelle \& Swain, 1970)}
\textbf{Variables:} $y_j\in\{0,1\}$ (abrir el sitio $j$); $x_{ij}\in[0,1]$ (asignar el cliente
$i$ al sitio $j$).
\begin{align}
  \min\ & \sum_{i=1}^{N}\sum_{j=1}^{M} d_{ij}\,x_{ij} \label{eq:f1obj}\\
  \text{s.a.}\ & \sum_{j=1}^{M} y_j = p \label{eq:f1p}\\
  & \sum_{j=1}^{M} x_{ij}=1 && i\in[N] \label{eq:f1asig}\\
  & x_{ij}\le y_j && i\in[N],\,j\in[M] \label{eq:f1link}\\
  & x_{ij}\ge 0,\quad y_j\in\{0,1\}. \nonumber
\end{align}
\textbf{Rol de cada bloque:} \eqref{eq:f1obj} suma distancias de las asignaciones activas;
\eqref{eq:f1p} abre exactamente $p$ sitios (restricción \emph{global}); \eqref{eq:f1asig} asigna
cada cliente a un sitio; \eqref{eq:f1link} \emph{acopla} $x$ con $y$ (solo se asigna a sitios
abiertos). Las $x$ pueden relajarse a continuas: con $y$ fijo y minimizando, cada cliente
concentra su peso en el sitio abierto más cercano, de modo que la solución es automáticamente
$0/1$.

\emph{Debilidad:} F1 tiene $N\times M$ variables $x$ y $1+N+N\times M$ restricciones; para
$N=M=10^4$ son $\sim 10^8$ variables, inviable. Esto motiva las formulaciones ordenadas.

\subsection{Distancias ordenadas y radios}
Para cada cliente $i$ sea $K_i\le M$ el número de \emph{distancias distintas} a los sitios,
ordenadas $D^1_i<D^2_i<\dots<D^{K_i}_i$. Se introduce $z^k_i=1$ para
$k=1,\dots,K_i-1$ si y solo si el cliente $i$ \emph{aún no está cubierto} al radio $D^k_i$ (no
hay sitio abierto a distancia $\le D^k_i$). No se necesita variable $z_i^{K_i}$: al último radio
ya se alcanzaron todos los sitios y, como $p\ge1$, algún sitio abierto existe.

\subsection{Por qué F2 mejora a F1: radios y umbral de distancia (Cornuejols, Nemhauser \& Wolsey, 1980)}
\textbf{La idea (antes del álgebra).} F1 razona por \emph{asignación}: ``¿a qué sitio mando cada
cliente?'', lo que cuesta una variable $x_{ij}$ por cada par. F2 razona por \emph{umbral de
distancia}: para cada cliente, ``¿a qué \textbf{radio} llega su sitio abierto más cercano?''.
Como muchos pares cliente--sitio comparten distancia (sobre todo si $N=M$), basta una variable
$z^k_i$ por cada \emph{radio intermedio} $D^k_i$, $k=1,\dots,K_i-1$, y los radios distintos
$K_i$ suelen ser muchos menos que $M$. Esa es la razón estructural por la que F2 nunca es mayor
que F1 y normalmente es mucho menor, manteniendo la misma cota LP.

\begin{align}
  \min\ & \sum_{i=1}^{N}\Big[\,D^1_i + \sum_{k=1}^{K_i-1}\big(D^{k+1}_i-D^k_i\big)z^k_i\,\Big] \label{eq:f2obj}\\
  \text{s.a.}\ & \sum_{j=1}^{M} y_j=p \nonumber\\
  & z^k_i + \!\!\sum_{j:\,d_{ij}\le D^k_i}\!\! y_j \ge 1 && i\in[N],\,k=1,\dots,K_i-1 \label{eq:f2cov}\\
  & z^k_i\ge 0 && i\in[N],\,k=1,\dots,K_i-1 \nonumber\\
  & y_j\in\{0,1\}. \nonumber
\end{align}

\begin{intuicion}[Peajes crecientes]
Un cliente con distancias posibles $2,5,9$ tiene objetivo $2+3\,z^1+4\,z^2$. Si su sitio abierto
más cercano está a $5$: ¿hay sitio dentro de radio $2$? No $\Rightarrow z^1=1$. ¿Dentro de $5$? Sí
$\Rightarrow z^2=0$. Costo $=2+3+0=5$. Se ``pagan peajes'' $D^{k+1}-D^k$ mientras no aparezca un
sitio abierto; al aparecer, se deja de pagar.
\end{intuicion}

La restricción \eqref{eq:f2cov} fuerza $z^k_i=1$ cuando no hay sitio abierto dentro de $D^k_i$.
F2 tiene $K=\sum_i (K_i-1)$ variables $z$ (con $K\le N\times M$, usualmente $K\ll N\times M$) y
\textbf{la misma relajación lineal} que F1.

\subsection{Por qué F3 mejora a F2: la misma formulación con matriz mucho más rala (Elloumi, 2010)}
Como los $z^k_i$ son monótonos decrecientes en $k$ ($z^{k-1}_i=0\Rightarrow z^k_i=0$), la
cobertura ``$\le D^k_i$'' puede reescribirse usando solo los sitios a distancia \emph{exactamente}
$D^k_i$:
\begin{align}
  \min\ & \sum_{i=1}^{N}\Big[\,D^1_i + \sum_{k=1}^{K_i-1}\big(D^{k+1}_i-D^k_i\big)z^k_i\,\Big] \label{eq:f3obj}\\
  \text{s.a.}\ & \sum_{j=1}^{M} y_j=p \nonumber\\
  & z^1_i + \!\!\sum_{j:\,d_{ij}=D^1_i}\!\! y_j \ge 1 && i\in[N] \label{eq:f3base}\\
  & z^k_i + \!\!\sum_{j:\,d_{ij}=D^k_i}\!\! y_j \ge z^{k-1}_i && i\in[N],\,k=2,\dots,K_i-1 \label{eq:f3rec}\\
  & z^k_i\ge 0 && i\in[N],\,k=1,\dots,K_i-1 \nonumber\\
  & y_j\in\{0,1\}. \nonumber
\end{align}
Al usar ``$=$'' en lugar de ``$\le$'', cada $y_j$ aparece en \emph{muchas menos} restricciones: la
matriz es mucho más rala. F2 y F3 comparten variables, objetivo y relajación lineal, pero F3
rinde \textbf{significativamente mejor} en la práctica por su \emph{esparsidad}. \textbf{Por eso
F3 es la base de la descomposición de Benders.}

\subsection{F4 — formulación compacta de cortes}
Escribir de una vez todos los cortes de Benders (Sección~\ref{sec:benders}) sobre una variable
$\theta_i$ por cliente da la forma compacta:
\begin{align}
  \min\ & \sum_{i=1}^N \theta_i \qquad\text{s.a.}\quad \sum_{j=1}^M y_j=p, \nonumber\\
  & \theta_i \ge D^1_i && i\in[N] \label{eq:f4base}\\
  & \theta_i \ge D^{k+1}_i-\!\!\sum_{j:\,d_{ij}\le D^k_i}\!\!\big(D^{k+1}_i-d_{ij}\big)y_j && i\in[N],\,k\in[K_i-1] \label{eq:f4cut}\\
  & y_j\in\{0,1\}. \nonumber
\end{align}
F4 tiene solo $N+M$ variables, pero \emph{tantas restricciones como F2} y matriz igual de densa.
Empíricamente, pasarle F4 \emph{entera} al solver es más lento que F2/F3: la densidad mata el
rendimiento. La clave no está en la formulación compacta, sino en \textbf{generar los cortes
\eqref{eq:f4cut} bajo demanda como \textit{lazy constraints}} dentro del branch-and-cut.

\subsection{Comparación estructural}
\begin{table}[h]\centering
\begin{tabular}{lcccc}
\toprule
 & \textbf{F1} & \textbf{F2} & \textbf{F3} & \textbf{F4} \\
\midrule
Variables (además de $y$) & $N\!\times\!M$ ($x$) & $K$ ($z$) & $K$ ($z$) & $N$ ($\theta$) \\
N.º restricciones & $1+N+NM$ & $1+K$ & $1+K$ & $1+K$ \\
Densidad de la matriz & media & alta & \textbf{baja} & alta \\
Relajación lineal & \multicolumn{4}{c}{idéntica en las cuatro} \\
Rol en este trabajo & referencia & previa & \textbf{base de Benders} & cortes generados \\
\bottomrule
\end{tabular}
\caption{$K=\sum_i (K_i-1)\le N\times M$. Las cuatro tienen el mismo óptimo entero y la misma cota LP;
difieren en cómo el solver las maneja.}
\end{table}
